\documentclass[10pt,a4paper]{article}
\usepackage[a4paper,margin=14mm,top=12mm,bottom=14mm]{geometry}
\usepackage{fontspec}
\setmainfont{Liberation Sans}
\usepackage{polyglossia}
\setdefaultlanguage{french}
\usepackage{amsmath}
\usepackage{xcolor}
\usepackage{graphicx}
\usepackage[most]{tcolorbox}
\usepackage{tabularx,booktabs,array,multirow}
\usepackage{enumitem}
\usepackage{fancyhdr}
\usepackage{lastpage}
\usepackage[hidelinks]{hyperref}
\definecolor{uimm}{HTML}{005677}
\definecolor{cyan}{HTML}{13A2DD}
\definecolor{fond}{HTML}{E6F2FA}
\definecolor{ok}{HTML}{2E7D32}
\definecolor{ko}{HTML}{B3261E}
\definecolor{violet}{HTML}{7B1FA2}
\setlength{\parindent}{0pt}
\setlength{\parskip}{3pt}
\setlist{nosep}
\setlist[itemize]{leftmargin=13pt,label=\textcolor{cyan}{\textbullet}}
\newcommand{\titre}[1]{\vspace{5pt}{\color{uimm}\large\bfseries #1}\par\vspace{-3pt}{\color{cyan}\rule{\linewidth}{1.2pt}}\par\vspace{0pt}}
\tcbset{boite/.style={colback=fond,colframe=uimm,boxrule=0.6pt,arc=2pt,left=6pt,right=6pt,top=4pt,bottom=4pt},
        defi/.style={colback=violet!6,colframe=violet,boxrule=0.8pt,arc=2pt,left=6pt,right=6pt,top=4pt,bottom=4pt,
                     title={\bfseries Défi},coltitle=white,colbacktitle=violet,fonttitle=\small}}
\newcommand{\bandeau}[3]{%
\begin{tcolorbox}[colback=uimm,colframe=uimm,arc=2pt,left=4pt,right=8pt,top=4pt,bottom=4pt]
\begin{minipage}[c]{18mm}\setlength{\fboxsep}{1.2pt}\fcolorbox{white}{white}{\includegraphics[height=16.5mm]{logo_uimm.jpg}}\end{minipage}\hspace{4mm}%
\begin{minipage}[c]{\dimexpr\linewidth-22mm\relax}\color{white}
{\small #1}\\[1pt]{\LARGE\bfseries #2}\\[2pt]{\large #3}\hfill{\small CQPM Technicien qualité}
\end{minipage}
\end{tcolorbox}}
\newcommand{\reponse}[1][1]{\par\vspace{2pt}{\color{gray!60}\foreach \i in {1,...,#1}{\rule{\linewidth}{0.3pt}\par\vspace{9pt}}}}
\usepackage{pgffor}
\pagestyle{fancy}\fancyhf{}
\renewcommand{\headrulewidth}{0pt}
\newcommand{\pied}[1]{\fancyfoot[L]{\footnotesize\color{uimm}#1}\fancyfoot[R]{\footnotesize\color{uimm}Page \thepage{} sur \pageref{LastPage}}}

\usepackage{tikz}
\usepackage{longtable}
\renewcommand{\bandeau}[3]{%
\begin{tcolorbox}[colback=uimm,colframe=uimm,arc=2pt,left=4pt,right=8pt,top=4pt,bottom=4pt]
\begin{minipage}[c]{18mm}\setlength{\fboxsep}{1.2pt}\fcolorbox{white}{white}{\includegraphics[height=16.5mm]{logo_uimm.jpg}}\end{minipage}\hspace{4mm}%
\begin{minipage}[c]{\dimexpr\linewidth-22mm\relax}\color{white}
{\small #1}\\[1pt]{\LARGE\bfseries #2}\\[2pt]{\large #3}
\end{minipage}
\end{tcolorbox}}
% en-tête d'une fiche : code, titre, question, bloc, durée, organisation, niveau
\newcommand{\entete}[7]{%
\begin{tcolorbox}[colback=uimm,colframe=uimm,arc=2pt,left=4pt,right=6pt,top=4pt,bottom=4pt]
\begin{minipage}[c]{18mm}\setlength{\fboxsep}{1.2pt}\fcolorbox{white}{white}{\includegraphics[height=16.5mm]{logo_uimm.jpg}}\end{minipage}\hspace{3mm}%
\begin{minipage}[c]{\dimexpr\linewidth-58mm\relax}\color{white}\raggedright
{\small Base SPC · #4}\\[1pt]{\Large\bfseries #1 · #2}\\[2pt]{\itshape #3}
\end{minipage}\hfill
\begin{minipage}[c]{33mm}\color{white}\small\raggedleft
Durée : \textbf{#5}\\ \textbf{#6}\\ Niveau #7 sur 3
\end{minipage}
\end{tcolorbox}}
\newcommand{\tagE}{{\setlength{\fboxsep}{1.5pt}\colorbox{cyan}{\color{white}\scriptsize\bfseries Ellistat}}}
\newcommand{\tagW}{{\setlength{\fboxsep}{1.5pt}\colorbox{uimm}{\color{white}\scriptsize\bfseries Page web}}}
\tcbset{pred/.style={colback=orange!6,colframe=orange!80!black,boxrule=0.6pt,arc=2pt,left=6pt,right=6pt,top=3pt,bottom=1pt,
                    title={\bfseries Je prédis, avant de commencer},coltitle=white,colbacktitle=orange!80!black,fonttitle=\small},
        retenir/.style={colback=fond,colframe=uimm,boxrule=0.6pt,arc=2pt,left=6pt,right=6pt,top=3pt,bottom=3pt,
                    title={\bfseries À retenir},coltitle=white,colbacktitle=uimm,fonttitle=\small},
        plus/.style={colback=violet!6,colframe=violet,boxrule=0.6pt,arc=2pt,left=6pt,right=6pt,top=3pt,bottom=3pt,
                    title={\bfseries Pour aller plus loin},coltitle=white,colbacktitle=violet,fonttitle=\small},
        act/.style={colback=uimm!10,colframe=uimm,boxrule=0.5pt,arc=2pt,left=5pt,right=5pt,top=2pt,bottom=2pt,before skip=7pt,after skip=3pt},
        verif/.style={colback=orange!6,colframe=orange!80!black,boxrule=0.6pt,arc=2pt,left=6pt,right=6pt,top=3pt,bottom=3pt,
                    title={\bfseries À vérifier sur poste avant la première séance},coltitle=white,colbacktitle=orange!80!black,fonttitle=\small}}
\newcommand{\fichier}[1]{\texttt{\small #1}}
\newcommand{\menu}[1]{\textsf{\textbf{#1}}}
\newcommand{\champ}[1]{\rule{#1}{0.3pt}}

\pied{Base SPC · fiche M16 · Deux grades, un micromètre}
\begin{document}
\thispagestyle{fancy}
\entete{M16}{Deux grades, un micromètre}{Mon micromètre voit-il la différence entre une cale de grade 0 et une cale de grade 1 ?}{Le moyen de mesure}{30 min}{en binôme}{2}
\begin{tcolorbox}[boite]\textbf{Le but.} Mesurer deux cales de même cote et de grades différents, et voir ce que le micromètre en dit.\par\textbf{Il vous faut.} Une cale céramique de 10 mm, grade 0 · Une cale acier de 10 mm, grade 1 · Un micromètre 0-25 mm · Ellistat\end{tcolorbox}
\begin{tcolorbox}[pred]La cale de grade 0 est plus précise que celle de grade 1. Votre micromètre au micron le verra-t-il ? Écrivez l'écart que vous attendez entre les deux moyennes, en µm.\reponse\end{tcolorbox}
\titre{À faire}
\begin{enumerate}[label=\textcolor{cyan}{\textbf{\arabic*.}},leftmargin=15pt,itemsep=1.5pt]
\item Lire dans le tableau ci-dessous les écarts limites des deux grades pour une cale de 10 mm. Calculer l'écart maximal possible entre les deux cales : les deux limites s'ajoutent.
\item Nettoyer les deux cales. Mesurer \textbf{15 fois} chacune au cliquet, en alternant : céramique, acier, céramique, acier. Retirer la cale à chaque mesure.
\item \tagE\ \menu{Gage R\&R}, onglet \menu{Justesse} : deux études, « Ceramique\_grade0 » et « Acier\_grade1 ». Nombre de mesures 15 ; Valeur vraie 10,000 ; LSI 9,925 et LSS 10,075 ; k 3. Relever Cg et Cgk.
\item \tagW\ Coller les deux séries dans la page « Comparer deux séries de mesures », avec 10,000 comme valeur de la cale.
\item Avec 15 mesures par cale, le plus petit écart de moyennes que le test déclare significatif vaut environ $0{,}75 \times s$. Prendre pour $s$ la moyenne des deux écarts-types.
\end{enumerate}
\vspace{3pt}{\renewcommand{\arraystretch}{1.35}\small\begin{tabularx}{\linewidth}{|l|*{3}{>{\centering\arraybackslash}X|}}\hline
\textbf{Écart limite de longueur (µm)} & \textbf{Grade 0} & \textbf{Grade 1} & \textbf{Grade 2}\\\hline
Cale jusqu'à 10 mm & $\pm 0{,}12$ & $\pm 0{,}20$ & $\pm 0{,}45$\\\hline
Cale de plus de 10 mm à 25 mm & $\pm 0{,}14$ & $\pm 0{,}30$ & $\pm 0{,}60$\\\hline
\end{tabularx}}\par{\footnotesize ISO 3650, valeurs relevées sur des pages de revendeurs : le certificat de vos cales fait foi.}\par\vspace{5pt}
\begin{minipage}[t]{0.485\linewidth}\textbf{Cale céramique grade 0 (mm)}\par\vspace{2pt}
{\renewcommand{\arraystretch}{1.4}\begin{tabularx}{\linewidth}{|*{3}{>{\centering\arraybackslash}p{5mm}|>{\centering\arraybackslash}X|}}\hline
\scriptsize 1 &  & \scriptsize 2 &  & \scriptsize 3 & \\\hline
\scriptsize 4 &  & \scriptsize 5 &  & \scriptsize 6 & \\\hline
\scriptsize 7 &  & \scriptsize 8 &  & \scriptsize 9 & \\\hline
\scriptsize 10 &  & \scriptsize 11 &  & \scriptsize 12 & \\\hline
\scriptsize 13 &  & \scriptsize 14 &  & \scriptsize 15 & \\\hline
\end{tabularx}}\end{minipage}\hfill\begin{minipage}[t]{0.485\linewidth}\textbf{Cale acier grade 1 (mm)}\par\vspace{2pt}
{\renewcommand{\arraystretch}{1.4}\begin{tabularx}{\linewidth}{|*{3}{>{\centering\arraybackslash}p{5mm}|>{\centering\arraybackslash}X|}}\hline
\scriptsize 1 &  & \scriptsize 2 &  & \scriptsize 3 & \\\hline
\scriptsize 4 &  & \scriptsize 5 &  & \scriptsize 6 & \\\hline
\scriptsize 7 &  & \scriptsize 8 &  & \scriptsize 9 & \\\hline
\scriptsize 10 &  & \scriptsize 11 &  & \scriptsize 12 & \\\hline
\scriptsize 13 &  & \scriptsize 14 &  & \scriptsize 15 & \\\hline
\end{tabularx}}\end{minipage}\par
\newpage
\titre{J'observe, je décide}
\begin{enumerate}[label=\textcolor{cyan}{\textbf{Q\arabic*.}},leftmargin=22pt]
\item Écart maximal possible entre les deux cales : \champ{14mm} µm. Écart observé entre vos deux moyennes : \champ{14mm} µm.\par\vspace{5pt}
\item La page conclut-elle à une différence réelle entre les deux moyennes ? Cela prouve-t-il que les deux cales sont identiques ?\reponse[2]
\item Plus petit écart détectable ($0{,}75 \times s$) = \champ{14mm} µm. Est-il plus petit ou plus grand que l'écart maximal possible entre les deux cales ?\reponse[1]
\item À quoi sert alors une cale de grade 0 ? Avec quel type d'instrument son grade devient-il visible ?\reponse[2]
\end{enumerate}
\begin{tcolorbox}[retenir]Un étalon doit être nettement meilleur que l'instrument qu'il qualifie. Plus fin que nécessaire, son grade ne se voit pas.\end{tcolorbox}
\begin{tcolorbox}[plus]Refaites avec 50 mesures par cale : le plus petit écart détectable tombe à environ $0{,}4 \times s$. Frôle-t-il l'écart maximal possible ? Que faudrait-il contrôler avant de conclure ?\end{tcolorbox}
\vfill{\footnotesize\color{uimm}\textbf{Ressources de la base :} Comparer deux séries de mesures · Données de secours.}
\end{document}